\chapter{主动热信息探针}
\label{ch:ati}

本章要判定的量是：一次边界变化是在买参数信息，还是在推动品质。二者可以驱动同一阀门。判定失败的表现是把每一次喷蒸汽都称为认知，或在方差没有下降时仍提前停机。

\section{探针与加工的区别}

\paragraph{要判定的量}
一段执行器轨迹是否同时满足信息目标、损伤预算和结束后的回退。

\paragraph{继承的文献层与不继承的层}
对偶控制的思想是输入同时推动状态与辨识 \cite{feldbaum1965}。信息度量继承熵差 \cite{lindley1956}。不继承在每个控制周期求解任意波形。家电执行器使用有限探针库。

\begin{definition}[热信息探针]
探针是一段有限时长 \(\tau\le\tau_{\max}\) 的执行器轨迹 \(a\in\cA\)，满足三条：设计目标是增大 \eqref{eq:fisher} 的 \(\log\det\)，或等价地减小决策量的后验熵；对表面褐变、质量损失和热力进度的预测损伤低于预算 \(\eta\)；结束后执行器回到探针前的名义轨迹，除非安全层已经改写该轨迹。
\end{definition}

不满足第一条的功率变化是加工动作，归第 \ref{ch:safe} 章调度。日志记录该时刻哪一个函数产生了阀门指令。

本次对偶策略用时间窗代替这个定义：\([150,190)\) 秒且方差大于 0.25 时进入探针模式。它没有在库上比较 \(G-\lambda C\)。因此它是一条信息动作的时间表，不是定义中的最优探针。

\paragraph{本节不能推出的结论}
时间窗内出现了蒸汽，不能推出该次动作满足损伤预算。预算需要表面褐变和冷点进度的预测增量。核在探针期间继续积分这些量，但没有单独的预算否决，除了加热器削减和厚件方差否决。

\section{信息目标}

\paragraph{要判定的量}
候选探针之间比较的是哪一个标量，以及这个标量在本次日志里有没有被计算。

高斯近似下，决策状态的微分熵由协方差的行列式决定。对候选探针 \(a\)，
\begin{equation}
\label{eq:ati}
G(a)=\frac12\log\det\vect{\Sigma}_{z,k}
-\frac12\log\det\vect{\Sigma}_{z,k+\tau}(a).
\end{equation}
在无偏估计的 Cramér–Rao 界附近，\(G(a)\) 与 Fisher 信息的增量同阶。在线规格使用集合形式。Fisher 形式在线性化工作点上平行计算。二者排序不一致时以集合形式为准，并记录该次失配。本次两种形式都没有计算。

代价
\begin{equation}
\label{eq:cost}
C(a)=w_E E(a)+w_m\Delta m(a)+w_b\Delta\alpha_b(a)+w_\tau\tau(a).
\end{equation}
最优探针
\begin{equation}
\label{eq:astar}
a^\star=\argmax_{a\in\cA}\big(G(a)-\lambda C(a)\big),
\end{equation}
仅当 \(G(a^\star)\ge G_{\min}\) 且被动估计已经失效时执行。薄件的 \(G_{\min}\) 应更高，因为被动观测本来就更有效；厚件的 \(G_{\min}\) 更低。本次没有这组阈值，时间窗对三种厚度一视同仁。交叉表里薄件对偶没有优势，与“薄件不该为了信息去探针”同方向，但不是 \(G_{\min}\) 标定的结果。

\subsection{若把这一次乘子当成一个标量方差}

已运行的更新不是 \eqref{eq:ati}。它在更接近参照时把一个标量乘 \(0.35\)。先验 \(0.64\)，乘完正好是 \(0.224\)。若把这个标量假装成一个高斯分量的方差，微分熵只差
\begin{equation}
\label{eq:nats}
\frac12\ln\frac{0.64}{0.224}
=\frac12\ln\frac{1}{0.35}
=0.525\ \mathrm{nat}
=0.757\ \mathrm{bit}.
\end{equation}
行列式在这里退化成一个正数。候选探针之间没有比较，所以这个差不能给探针排序。同一乘子打在每一次被接受的更新上，信息目标对探针库是常数。

同一标量还留在加热裕量里。厚度 \(30\,\mathrm{mm}\) 时，更新前
\[
m_g=6+10\sqrt{0.64}=14,\qquad 75+m_g=89^\circ\mathrm{C}.
\]
更新后 \(\sqrt{0.224}=0.473\)，裕量 \(10.73\,\mathrm{K}\)，阈值 \(85.73^\circ\mathrm{C}\)。\(\lambda_I\) 没有被置零。加热律仍带着新的裕量。种子 7 在 \(192\,\mathrm{s}\) 之后方差保持 \(0.224\)，所以后半程的阈值停在 \(85.73^\circ\mathrm{C}\)，不是回到 \(75^\circ\mathrm{C}\)。

\paragraph{本节不能推出的结论}
不能把方差从 0.64 乘到 0.224 写成 \(G(a)\) 的奈特数。\eqref{eq:nats} 是把乘子假装成标量高斯之后的算术，不是核算出的信息增益。

\section{探针库，而不是连续最优控制}

\paragraph{要判定的量}
库里每一条轨迹主要敏感的参数是哪一个，以及单独蒸汽为什么不够。

\begin{table}[ht]
\centering
\caption{探针库}
\begin{tabularx}{\textwidth}{lXll}
\toprule
编号 & 轨迹 & 主要敏感参数 & 时长 \\
\midrule
P1 & 蒸汽开、风机低、加热保持 & 冷凝系数、表面热容 & 4--8 s \\
P2 & 蒸汽关、排汽开、风机高 & 外部热阻、传质系数 & 6--10 s \\
P3 & 加热阶跃、风机关 & 辐射份额、发射率路径 & 8--12 s \\
P4 & P1 之后立即 P2 & 把冷凝跃迁与对流热阻分开 & \(\le 20\) s \\
P0 & 维持现状 & \(G\) 的零点 & 0 \\
\bottomrule
\end{tabularx}
\end{table}

P4 是规格中的默认信息探针。单独的 P1 把冷凝潜热和内部导热缠在一起。紧接着的排汽和高风只改变外部热阻与蒸发方向。两段响应之差对内部热阻敏感。命题 \ref{prop:pulse} 把这件事写成秩：第二段对内部热阻的特解不能写成第一段的标量倍。

断言用两节点比较了“蒸汽然后高风”和“只喷蒸汽”。滞后分别为 40.846 和 38.912。判据是差值大于 1，不是大于 15\%。15\% 的相对判据在这对数字上不成立，因为 \(40.846/38.912\approx 1.050\)。正文只使用已经通过的绝对差。

已运行的对偶窗是 40 s，长于表中 P4 的 20 s 上限。170 s 之前蒸汽指令为 1，其后风机为 1.3。加热请求在窗内仍由冷点跟踪器决定，不是“加热保持”的常数。所以这个窗只保留了 P4 的顺序——先冷凝方向、后对流——没有保留库的时长和打分。

\begin{figure}[ht]
\centering
\begin{tikzpicture}[x=0.18cm,y=0.7cm]
\draw[->] (0,0) -- (46,0) node[right] {\(t/\mathrm{s}\)};
\draw (15,0) -- (15,1.2) -- (17,1.2) -- (17,0);
\draw (17,0) -- (17,0.7) -- (19,0.7) -- (19,0);
\node[above] at (16,1.2) {蒸汽};
\node[above] at (18,0.7) {高风};
\node at (8,0.4) {P0};
\node at (30,0.4) {回到名义轨迹};
\end{tikzpicture}
\caption{P4 两段时序的依存关系图，不是某次日志的执行器曲线。已运行窗是 150 s 至 190 s，长于图中的规格上限。}
\end{figure}

\paragraph{本节不能推出的结论}
库有五个元素，不能推出五个元素都在 S1 里跑过。S1 的对偶策略只有固定窗。P0、P1、P2、P3 的信息增益比较没有日志。

\section{损伤预算如何算}

\paragraph{要判定的量}
短探针的损伤写在冷点上，还是写在表面上。

用 \eqref{eq:kinetics} 在表面节点上积分褐变增量，用 \eqref{eq:F} 在冷点上确认热力增量。冷点增量的上界由命题 \ref{prop:budget} 给出：窗口 \(\tau\le\tau_2/10\) 且表面增量不超过 \(\Delta T_s\) 时，
\begin{equation}
|\Delta T_2|\le \frac{\tau}{\tau_2}\Delta T_s.
\end{equation}
断言取 \(\tau_2=30\)、\(\tau=3\)、\(\Delta T_s=12\)，上界 1.2，离散积分得到 1.143。上界成立。它使用的时间常数比本次 40 s 窗口更严格。把 40 s 代进 30 mm 的导热尺度，Fourier 数仍远小于 1，但那是另一组常数，不能把 1.143 说成 S1 探针的冷点增量。

种子 7 在 152 s 至 192 s 的冷点增量约为 \(0.04\,\mathrm{K}\)。表面增量约为 \(8.5\,\mathrm{K}\)。这一条轨迹上，损伤确实不在冷点。表面在同一窗口已经是 \(80^\circ\mathrm{C}\) 至 \(89^\circ\mathrm{C}\)，褐变公式要求表面超过 \(110^\circ\mathrm{C}\) 且含水低于 0.55 才增长，所以这一窗的褐变增量也为零。不能把这一条推广成“探针总是无害”：更早的加热已经把表面送到 \(80^\circ\mathrm{C}\)，探针是叠加上去的。

硬否决不进入 \eqref{eq:astar}：水位低于干烧阈值，门打开，腔压或门封温度接近安全边界，表面温度已在褐变失控区时再做加热阶跃。这些测量本次没有采样，否决不存在于步进里。

\paragraph{本节不能推出的结论}
冷点上界不否决表面已经越界的探针。用冷点温度决定“探针无害”会放过表面。

\section{在线算法}

规格中的在线顺序是六步。每个估计周期计算信息赤字。赤字消失且进度的下分位已经高于设计线时，不探针。否则对库中五条轨迹向前滚，计算 \(G\) 与 \(C\)，滚动时写入执行器饱和与传感器饱和。选取最优；若最优仍是维持现状，则保持加工轨迹并提高不确定度提示。执行期间冻结慢速菜单，快安全环继续。结束后若协方差没有降到预测的 70\% 以内，标记失配并切换到更保守的约束。

第 6 步防止动作做了而信息没有进来。本次实现是这条原则的一个特例，不是六步算法：窗只在方差大于 0.25 时进入；更新只在更接近参照时把方差乘 0.35；厚件上方差仍大于 0.45 时不准停止。没有向前滚动，没有 70\% 的预测对照。种子 7 的方差在 192 s 变为 \(0.224\)，比值正好是 0.35，说明走的是乘子，不是 70\% 检验。

\paragraph{本节不能推出的结论}
不能把“标记失配”的次数写成结果。没有这个计数器。

\section{对偶目标}

探针库解决的是执行器不能任意波形。决策层仍把信息与品质放在同一次优化里，避免两个状态机抢执行器。
\begin{equation}
\label{eq:dual}
\begin{aligned}
\max_{\vect{u}}\;
& \E\big[Q(\vect{z}_{k+H})\big]
+\lambda_I\, G(\vect{u})
-\lambda_E E(\vect{u})\\
\mathrm{s.t.}\quad
& \Delta\alpha_b\le\eta,\quad
\Delta m\le m_{\mathrm{budget}},\quad
\vect{u}\in\cU.
\end{aligned}
\end{equation}
\(\lambda_I\) 随信息赤字增大。赤字低于零且安全分位已经通过时，\(\lambda_I=0\)，最优解退回纯品质控制。实现时先在库上比较，再把胜出的探针作为预测控制时域中的一段固定输入。

本次没有这个最大化。对偶与无探针两支共用同一个饱和线性加热律。差别只在 40 s 的窗和其后的热阻更新，以及厚件否决。能量中位数从无探针的 1655 到对偶的 1730，差在这一个窗和稍晚的停机上，不是 \(\lambda_E\) 标定的结果。

\paragraph{本节不能推出的结论}
\(\lambda_I=0\) 的极限没有在日志里单独作为一组策略跑出。对偶策略在方差已经下降之后仍使用带裕量的加热律，裕量里还留着 \(\sqrt{\mathrm{Var}}\)。信息项并没有在赤字下降后被关掉。

\section{短脉冲几乎不推动冷点}

表面时间常数 \(\tau_1=R_{\mathrm{ext}}C_1\)，内部时间常数 \(\tau_2=R_{\mathrm{int}}C_2\)。在 \(\tau\ll\tau_2\) 的窗口内，冷点温度增量是 \(O(\tau/\tau_2)\)。进度增量同阶。厚度增大使窗口更宽，同时使表面响应对内部热阻更不敏感。信息增益因此有上界。上界低于停机所需的熵减时，界面输出区间，不输出一个剩余时间点。

\begin{figure}[ht]
\centering
\begin{tikzpicture}[x=1.1cm,y=1cm]
\draw[->] (0,0) -- (6.2,0) node[right] {时间};
\draw[->] (0,0) -- (0,3.2) node[above] {响应};
\draw[domain=0:1.2,smooth] plot (\x,{2.4*(1-exp(-3*\x))});
\draw[domain=0:5.5,smooth] plot (\x,{2.4*(1-exp(-0.35*\x))});
\draw[dashed] (1.2,0) -- (1.2,2.6);
\node[below] at (1.2,0) {\(\tau\)};
\node[right] at (2.2,2.2) {表面};
\node[right] at (3.5,1.2) {冷点};
\end{tikzpicture}
\caption{表面时间常数与冷点窗口的依存关系图，不是测量。短窗内冷点曲线几乎尚未离开初值。}
\end{figure}

仿真规格要求在厚度与含水平面上计算这个上界。本次没有这张等高线。附录 \ref{app:scales} 只算了 Fourier 数和 Biot 数。种子 7 提供一条同方向的抽样，不能填这张图。

\paragraph{本节不能推出的结论}
“几乎不推动”不能写成任意时长的蒸汽都安全。时长一旦接近内部时间常数，上界的前提 \(\tau\le\tau_2/10\) 不再成立。
